\documentclass[12pt, a4paper]{article}
\usepackage[UTF8]{ctex}
\usepackage{amsmath}
\usepackage{amssymb}
\usepackage{geometry}
\usepackage{enumitem}
\usepackage{fancyhdr}
\usepackage{titlesec}

% 页面设置
\geometry{left=2.5cm, right=2.5cm, top=2.5cm, bottom=2.5cm}

% 页眉页脚
\pagestyle{fancy}
\fancyhf{}
\rhead{专升本数据结构习题集}
\lhead{分类整理}
\cfoot{\thepage}

% 标题格式
\titleformat{\section}{\large\bfseries}{\thesection}{1em}{}
\titleformat{\subsection}{\normalsize\bfseries}{\thesubsection}{1em}{}

\title{\textbf{第六章 树}}
\author{}
\date{}

\begin{document}

\section*{选择题}

\begin{enumerate}[label=\arabic*.]
	\item 树最适合用来表示（ ）
	      \begin{enumerate}
		      \item[A.] 有序数据元素
		      \item[B.] 无序数据元素
		      \item[C.] 元素之间具有分支层次关系的数据
		      \item[D.] 元素之间无联系的数据
	      \end{enumerate}

	\item 在一棵树中，（ ）没有前驱结点。
	      \begin{enumerate}
		      \item[A.] 分支结点
		      \item[B.] 叶结点
		      \item[C.] 树根结点
		      \item[D.] 空结点
	      \end{enumerate}

	\item 有 $N$ 个叶子结点的哈夫曼树中，其结点总数为（ ）
	      \begin{enumerate}
		      \item[A.] 不确定
		      \item[B.] $2N$
		      \item[C.] $2N+1$
		      \item[D.] $2N-1$
	      \end{enumerate}

	\item 把一棵树转换成对应的二叉树时，这棵二叉树的形态是（ ）
	      \begin{enumerate}
		      \item[A.] 唯一的
		      \item[B.] 有多种
		      \item[C.] 有多种但根结点都没有左子树
		      \item[D.] 有多种但根结点都没有右子树
	      \end{enumerate}

	\item 树的后根遍历序列与其对应的二叉树的后序遍历序列相同。（ ）
	      \begin{enumerate}
		      \item[A.] 正确
		      \item[B.] 错误
	      \end{enumerate}

	\item 树的先根遍历序列与其对应的二叉树的中序遍历序列相同。（ ）
	      \begin{enumerate}
		      \item[A.] 正确
		      \item[B.] 错误
	      \end{enumerate}

	\item 按照二叉树的定义，具有 3 个结点的二叉树有（ ）种。
	      \begin{enumerate}
		      \item[A.] 3
		      \item[B.] 4
		      \item[C.] 5
		      \item[D.] 6
	      \end{enumerate}

	\item 深度为 5 的二叉树至多有（ ）个结点。
	      \begin{enumerate}
		      \item[A.] 16
		      \item[B.] 32
		      \item[C.] 31
		      \item[D.] 10
	      \end{enumerate}

	\item 在一棵非空二叉树的中序遍历序列中，根结点的右边（ ）
	      \begin{enumerate}
		      \item[A.] 只有右子树上的所有结点
		      \item[B.] 只有右子树上的部分结点
		      \item[C.] 只有左子树上的部分结点
		      \item[D.] 只有左子树上的所有结点
	      \end{enumerate}

	\item 任何一棵二叉树的叶子结点在先序，中序和后序遍历序列中的相对次序（ ）
	      \begin{enumerate}
		      \item[A.] 不发生改变
		      \item[B.] 发生改变
		      \item[C.] 不能确定
	      \end{enumerate}

	\item 对一个满二叉树，m 个树叶，n 个结点，深度为 h，则（ ）
	      \begin{enumerate}
		      \item[A.] $n=h+m$
		      \item[B.] $h+m=2n$
		      \item[C.] $m=h-1$
		      \item[D.] $n=2^h-1$
	      \end{enumerate}

	\item 设 n, m 为一棵二叉树上的两个结点，在中序遍历时，n 在 m 前的条件是（ ）
	      \begin{enumerate}
		      \item[A.] n 在 m 右方
		      \item[B.] n 是 m 祖先
		      \item[C.] n 在 m 左方
		      \item[D.] n 是 m 子孙
	      \end{enumerate}

	\item 线索二叉树是一种（ ）结构。
	      \begin{enumerate}
		      \item[A.] 逻辑
		      \item[B.] 逻辑和存储
		      \item[C.] 物理
		      \item[D.] 线性结构
	      \end{enumerate}

	\item 在线索化树中，每个结点必须设置一个标志来说明它的左、右链指向的是树结构信息，还是线索化信息，若 0 标识树结构信息，1 标识线索，对应叶结点的左右链域，应标识为（ ）
	      \begin{enumerate}
		      \item[A.] 00
		      \item[B.] 01
		      \item[C.] 10
		      \item[D.] 11
	      \end{enumerate}

	\item 引入线索树的目的是（ ）
	      \begin{enumerate}
		      \item[A.] 加快查找结点的前驱或后继的速度
		      \item[B.] 为了能在二叉树中方便的进行插入与删除
		      \item[C.] 为了能方便的找到双亲
		      \item[D.] 使二叉树的遍历结果唯一
	      \end{enumerate}

	\item 根据使用频率，为 5 个字符设计的哈夫曼编码不可能的是（ ）
	      \begin{enumerate}
		      \item[A.] 111, 110, 10, 01, 00
		      \item[B.] 000, 001, 010, 011, 1
		      \item[C.] 100, 11, 10, 1, 0
		      \item[D.] 001, 000, 01, 11, 10
	      \end{enumerate}

	\item 根据使用频率，为 5 个字符设计的哈夫曼编码不可能的是（ ）
	      \begin{enumerate}
		      \item[A.] 000, 001, 010, 011, 1
		      \item[B.] 0000, 0001, 001, 01, 1
		      \item[C.] 000, 001, 01, 10, 11
		      \item[D.] 00, 100, 101, 110, 111
	      \end{enumerate}

	\item 已知一算术表达式的中缀形式为 $A+B*C-D/E$，后缀形式为 $ABC*+DE/-$，其前缀形式为（ ）
	      \begin{enumerate}
		      \item[A.] $-A+B*C/DE$
		      \item[B.] $-A+B*CD/E$
		      \item[C.] $-+*ABC/DE$
		      \item[D.] $-+A*BC/DE$
	      \end{enumerate}

	\item 设有一表示算术表达式的二叉树（见下图），它所表示的算术表达式是（ ）
	      \begin{enumerate}
		      \item[A.] $A*B+C/(D*E)+(F-G)$
		      \item[B.] $(A*B+C)/(D*E)+(F-G)$
		      \item[C.] $(A*B+C)/(D*E+(F-G))$
		      \item[D.] $A*B+C/D*E+F-G$
	      \end{enumerate}

	\item 在下述结论中，正确的是（ ）
	      \begin{enumerate}
		      \item[①] 只有一个结点的二叉树的度为 0;
		      \item[②] 二叉树的度为 2；
		      \item[③] 二叉树的左右子树可任意交换;
		      \item[④] 深度为 K 的完全二叉树的结点个数小于或等于深度相同的满二叉树。
		      \item[A.] ①②③
		      \item[B.] ②③④
		      \item[C.] ②④
		      \item[D.] ①④
	      \end{enumerate}

	\item 设森林 F 对应的二叉树为 B，它有 m 个结点，B 的根为 p，p 的右子树结点个数为 n，森林 F 中第一棵树的结点个数是（ ）
	      \begin{enumerate}
		      \item[A.] $m-n$
		      \item[B.] $m-n-1$
		      \item[C.] $n+1$
		      \item[D.] 条件不足，无法确定
	      \end{enumerate}

	\item 对于有 n 个结点的二叉树，其高度为（ ）
	      \begin{enumerate}
		      \item[A.] $n \log n$
		      \item[B.] $\log n$
		      \item[C.] $\lfloor \log n \rfloor + 1$
		      \item[D.] 不确定
	      \end{enumerate}

	\item 深度为 h 的满 m 叉树的第 k 层有（ ）个结点。($1 \le k \le h$)
	      \begin{enumerate}
		      \item[A.] $m^{k-1}$
		      \item[B.] $m^k-1$
		      \item[C.] $m^{h-1}$
		      \item[D.] $m^k$
	      \end{enumerate}

	\item 在下列存储形式中，哪一个不是树的存储形式？（ ）
	      \begin{enumerate}
		      \item[A.] 双亲表示法
		      \item[B.] 孩子链表表示法
		      \item[C.] 孩子兄弟表示法
		      \item[D.] 顺序存储表示法
	      \end{enumerate}

	\item 引入二叉线索树的目的是（ ）
	      \begin{enumerate}
		      \item[A.] 加快查找结点的前驱或后继的速度
		      \item[B.] 为了能在二叉树中方便的进行插入与删除
		      \item[C.] 为了能方便的找到双亲
		      \item[D.] 使二叉树的遍历结果唯一
	      \end{enumerate}

	\item 在叶子数目和权值相同的所有二叉树中，最优二叉树一定是完全二叉树，该说法（ ）
	      \begin{enumerate}
		      \item[A.] 正确
		      \item[B.] 错误
	      \end{enumerate}

	\item 某二叉树结点的中序序列为 ABCDEFG，后序序列为 BDCAFGE，则其左子树中结点数目为：（ ）
	      \begin{enumerate}
		      \item[A.] 3
		      \item[B.] 2
		      \item[C.] 4
		      \item[D.] 5
	      \end{enumerate}

	\item 某二叉树的前序遍历序列和中序遍历序列分别为 abc 和 bca，该二叉树的后序遍历序列是（ ）
	      \begin{enumerate}
		      \item[A.] cba
		      \item[B.] bca
		      \item[C.] abc
		      \item[D.] acb
	      \end{enumerate}

	\item 某二叉树的后序遍历序列和中序遍历序列分别为 cba 和 bca，该二叉树的前序遍历序列是（ ）
	      \begin{enumerate}
		      \item[A.] abc
		      \item[B.] cba
		      \item[C.] acb
		      \item[D.] bca
	      \end{enumerate}

	\item 若一棵二叉树具有 10 个度为 2 的结点，5 个度为 1 的结点，则度为 0 的结点的个数是（ ）
	      \begin{enumerate}
		      \item[A.] 9
		      \item[B.] 11
		      \item[C.] 15
		      \item[D.] 不能确定
	      \end{enumerate}

	\item 具有 10 个叶子结点的二叉树中有（ ）个度为 2 的结点。
	      \begin{enumerate}
		      \item[A.] 8
		      \item[B.] 9
		      \item[C.] 10
		      \item[D.] 11
	      \end{enumerate}

	\item 在具有 n 个度数为 2 的二叉树中，必有（ ）个叶子结点。
	      \begin{enumerate}
		      \item[A.] $n+2$
		      \item[B.] $n+1$
		      \item[C.] $n$
		      \item[D.] $n-1$
	      \end{enumerate}

	\item 假设二叉树对应的二叉链表共有 n 个非空链域，那么该二叉树有（ ）个结点的二叉树。
	      \begin{enumerate}
		      \item[A.] $n-1$
		      \item[B.] $n$
		      \item[C.] $n+1$
		      \item[D.] $2n$
	      \end{enumerate}

	\item n 个结点的二叉树，其对应的二叉链表共有（ ）个非空链域。
	      \begin{enumerate}
		      \item[A.] $n+1$
		      \item[B.] $n-1$
		      \item[C.] $2n$
		      \item[D.] $n$
	      \end{enumerate}

	\item m 个结点的二叉树，其对应的二叉链表共有（ ）个非空链域。
	      \begin{enumerate}
		      \item[A.] $m+1$
		      \item[B.] $2m$
		      \item[C.] $m-1$
		      \item[D.] $m$
	      \end{enumerate}

	\item 若二叉树对应的二叉链表共有 n 个非空链域，则该二叉树有（ ）个结点。
	      \begin{enumerate}
		      \item[A.] $n+1$
		      \item[B.] $n$
		      \item[C.] $n-1$
		      \item[D.] $2n$
	      \end{enumerate}

	\item 深度为 n 的完全二叉树最多有（ ）个结点。
	      \begin{enumerate}
		      \item[A.] $2n+1$
		      \item[B.] $2n-1$
		      \item[C.] $2^n$
		      \item[D.] $2^n-1$
	      \end{enumerate}

	\item 深度为 n 的完全二叉树最多有（ ）个结点。
	      \begin{enumerate}
		      \item[A.] $2^n$
		      \item[B.] $2^{n-1}$
		      \item[C.] $2^n-1$
		      \item[D.] $2^{n+1}-1$
	      \end{enumerate}

	\item 深度为 h 的二叉树，第 h 层至少有（ ）个结点。
	      \begin{enumerate}
		      \item[A.] 1
		      \item[B.] 0
		      \item[C.] 2
		      \item[D.] $2^{h-1}$
	      \end{enumerate}

	\item 深度为 h 的二叉树，第 h 层最多有（ ）个结点。
	      \begin{enumerate}
		      \item[A.] $h$
		      \item[B.] $2h$
		      \item[C.] $2^{h-1}$
		      \item[D.] $2^h$
	      \end{enumerate}

	\item 如果一个 huffman 树含有 n 个叶子，则该树必有（ ）的结点。
	      \begin{enumerate}
		      \item[A.] $2n$
		      \item[B.] $2n-1$
		      \item[C.] $n+1$
		      \item[D.] $n-1$
	      \end{enumerate}

	\item 一棵有 n 个叶子结点的哈夫曼树共有（ ）个结点。
	      \begin{enumerate}
		      \item[A.] $n$
		      \item[B.] $n-1$
		      \item[C.] $2n-1$
		      \item[D.] $2n$
	      \end{enumerate}

	\item 已知某棵二叉树的前序遍历序列是 abdgcefh，中序序列是 dgbaechf，其后序序列为（ ）
	      \begin{enumerate}
		      \item[A.] gdbehfca
		      \item[B.] gdbhefca
		      \item[C.] gdbehfac
		      \item[D.] gdbehfac
	      \end{enumerate}

	\item 线索二叉树的左线索指向其遍历序列中的（ ），右线索指向其遍历序列中的（ ）
	      \begin{enumerate}
		      \item[A.] 前驱 后继
		      \item[B.] 后继 前驱
	      \end{enumerate}

	\item 下面关于二叉树的叙述中，正确的是（ ）
	      \begin{enumerate}
		      \item[A.] 二叉树中至少存在一个度为 2 的结点
		      \item[B.] 二叉树中不存在度大于 2 的结点
		      \item[C.] 二叉树中不存在度为 1 的结点
		      \item[D.] 二叉树中只存在度为 0 的叶子结点
	      \end{enumerate}

	\item 在哈夫曼树中，权值最小的结点离根结点（ ）
	      \begin{enumerate}
		      \item[A.] 最近
		      \item[B.] 最远
		      \item[C.] 距离相等
		      \item[D.] 无法确定
	      \end{enumerate}

	\item 若已知一棵二叉树的前序遍历序列和后序遍历序列，则（ ）恢复该二叉树。
	      \begin{enumerate}
		      \item[A.] 可以
		      \item[B.] 不可以
	      \end{enumerate}

	\item 在二叉树中，指针 p 所指结点为叶子结点的条件是（ ）
	      \begin{enumerate}
		      \item[A.] p->lchild == NULL \&\& p->rchild == NULL
		      \item[B.] p->lchild == NULL || p->rchild == NULL
		      \item[C.] p->lchild == p->rchild
		      \item[D.] p->lchild + p->rchild == 0
	      \end{enumerate}

	\item 二叉树为二叉排序树的充分必要条件是其任一结点的值均大于其左孩子的值、小于其右孩子的值。这种说法（ ）
	      \begin{enumerate}
		      \item[A.] 正确
		      \item[B.] 错误
	      \end{enumerate}

	\item 下列叙述中，正确的是（ ）
	      \begin{enumerate}
		      \item[A.] 深度为 k 的二叉树最多有 $2^k$ 个结点
		      \item[B.] 深度为 k 的二叉树最多有 $2^{k-1}$ 个结点
		      \item[C.] 深度为 k 的二叉树最多有 $2^k - 1$ 个结点
		      \item[D.] 深度为 k 的二叉树最多有 $2^{k+1}$ 个结点
	      \end{enumerate}

	\item 若某二叉树有 20 个叶子结点，有 30 个结点仅有一个孩子，则该二叉树的总结点个数为（ ）
	      \begin{enumerate}
		      \item[A.] 49
		      \item[B.] 50
		      \item[C.] 69
		      \item[D.] 51
	      \end{enumerate}

	\item 在哈夫曼树中，权值最小的结点离根结点最近。（ ）
	      \begin{enumerate}
		      \item[A.] 正确
		      \item[B.] 错误
	      \end{enumerate}

	\item 强连通图的各顶点间均可达。（ ）
	      \begin{enumerate}
		      \item[A.] 正确
		      \item[B.] 错误
	      \end{enumerate}

	\item 对于任意一个图，从它的某个结点进行一次深度或广度优先遍历可以访问到该图的每个顶点。（ ）
	      \begin{enumerate}
		      \item[A.] 正确
		      \item[B.] 错误
	      \end{enumerate}

	\item 度为 2 的有序树是二叉树。（ ）
	      \begin{enumerate}
		      \item[A.] 正确
		      \item[B.] 错误
	      \end{enumerate}

	\item 二叉树的前序遍历序列中，任意一个结点均处在其孩子结点的前面。（ ）
	      \begin{enumerate}
		      \item[A.] 正确
		      \item[B.] 错误
	      \end{enumerate}

	\item 二叉树的后序遍历序列中，任意一个结点均处在其孩子结点的后面。（ ）
	      \begin{enumerate}
		      \item[A.] 正确
		      \item[B.] 错误
	      \end{enumerate}

	\item 用一维数组存储二叉树时，总是以前序遍历顺序存储结点。（ ）
	      \begin{enumerate}
		      \item[A.] 正确
		      \item[B.] 错误
	      \end{enumerate}

	\item 平衡二叉树中，任意结点左右子树的高度差（绝对值）不超过 1。（ ）
	      \begin{enumerate}
		      \item[A.] 正确
		      \item[B.] 错误
	      \end{enumerate}

	\item 在平衡二叉树中，任意结点左右子树的高度差（绝对值）不超过 1。（ ）
	      \begin{enumerate}
		      \item[A.] 正确
		      \item[B.] 错误
	      \end{enumerate}
\end{enumerate}

\section*{填空题}

\begin{enumerate}
	\item 二叉树由（ ），（ ），（ ）三个基本单元组成。
	\item 树在计算机内的表示方式有（ ），（ ），（ ）。
	\item 在二叉树中，指针 p 所指结点为叶子结点的条件是（ ）。
	\item 中缀式 $a+b*3+4*(c-d)$ 对应的前缀式为（ ），若 $a=1, b=2, c=3, d=4$，则后缀式 $db/cc*a-b*+$ 的运算结果为（ ）。
	\item 二叉树中某一结点左子树的深度减去右子树的深度称为该结点的（ ）。
	\item 具有 256 个结点的完全二叉树的深度为（ ）。
	\item 已知一棵度为 3 的树有 2 个度为 1 的结点，3 个度为 2 的结点，4 个度为 3 的结点，则该树有（ ）个叶子结点。
	\item 深度为 k 的完全二叉树至少有（ ）个结点，至多有（ ）个结点。
	\item 在完全二叉树中，编号为 i 和 j 的两个结点处于同一层的条件是（ ）。
	\item 假设根结点的层数为 1，具有 n 个结点的二叉树的最大高度是（ ）。
	\item 在一棵二叉树中，度为零的结点的个数为 $N_0$，度为 2 的结点的个数为 $N_2$，则有 $N_0=$（ ）。
	\item 设有 N 个结点的完全二叉树顺序存放在向量 A[1:N] 中，其下标值最大的分支结点为（ ）。
	\item 高度为 K 的完全二叉树至少有（ ）个叶子结点。
	\item 已知二叉树有 50 个叶子结点，则该二叉树的总结点数至少是（ ）。
	\item 如果结点 A 有 3 个兄弟，而且 B 是 A 的双亲，则 B 的度是（ ）。
	\item 对于一个具有 n 个结点的二叉树，当它为一棵（ ）二叉树时具有最小高度，当它为一棵（ ）时，具有最大高度。
	\item 具有 N 个结点的二叉树，采用二叉链表存储，共有（ ）个空链域。
	\item 含 4 个度为 2 的结点和 5 个叶子结点的二叉树，可有（ ）个度为 1 的结点。
	\item 二叉树结点的中序序列为 A, B, C, D, E, F, G，后序序列为 B, D, C, A, F, G, E，则该二叉树结点的前序序列为（ ），则该二叉树对应的树林包括（ ）棵树。
	\item 现有按中序遍历二叉树的结果为 abc，问有（ ）种不同的二叉树可以得到这一遍历结果，这些二叉树分别是（ ）。
	\item 若一个二叉树的叶子结点是某子树的中序遍历序列中的最后一个结点，则它必是该子树的（ ）序列中的最后一个结点。
	\item 先根次序周游树林正好等同于按（ ）周游对应的二叉树；后根次序周游树林正好等同于（ ）周游对应的二叉树。
	\item 具有 n 个结点的满二叉树，其叶子结点的个数是（ ），分支结点数是（ ）。
	\item 线索二叉树的左线索指向其（ ），右线索指向其（ ）。
	\item 哈夫曼树是（ ）。
	\item 若以 $\{4, 5, 6, 7, 8\}$ 作为叶子结点的权值构造哈夫曼树，则其带权路径长度是（ ）。
	\item 设 n 为哈夫曼树的叶子结点数目，则该哈夫曼树共有（ ）个结点。
	\item 在一棵二叉树中，度为零的结点的个数为 $n_0$，度为 2 的结点的个数为 $n_2$，则有 $n_0 =$ \underline{$n_2+1$}。
	\item 在有 n 个结点的二叉链表中，空链域的个数为 \underline{$n+1$}。
	\item 一棵有 n 个叶子结点的哈夫曼树共有 \underline{$2n-1$} 个结点。
	\item 深度为 5 的二叉树至多有 \underline{31} 个结点。
	\item 线索二叉树的左线索指向其遍历序列中的 \underline{前驱}，右线索指向其遍历序列中的 \underline{后继}。
	\item 在哈夫曼树中，权值最小的结点离根结点最远。
	\item 在平衡二叉树中，任意结点左右子树的高度差（绝对值）不超过 1。
	\item 若某二叉树有 20 个叶子结点，有 30 个结点仅有一个孩子，则该二叉树的总结点个数为 69。
\end{enumerate}

\section*{应用题}

\begin{enumerate}
	\item
	      \begin{enumerate}
		      \item[①] 试找出满足下列条件的二叉树：
		            1) 先序序列与后序序列相同
		            2) 中序序列与后序序列相同
		            3) 先序序列与中序序列相同
		            4) 中序序列与层次遍历序列相同
		            5) 先序序列与中序序列相反
		      \item[②] 已知一棵二叉树的中序序列和后序序列分别为 DBEAFIHCG 和 DEBHIFGCA，画出这棵二叉树，并求该二叉树的先序序列。
	      \end{enumerate}

	\item 将下列由三棵树组成的森林转换为二叉树。（只要求给出转换结果）

	\item 先序遍历序列：A B D F C E G H；中序遍历序列：B F D A G E H C
	      \begin{enumerate}
		      \item 画出这棵二叉树。
		      \item 画出这棵二叉树的后序线索树。
		      \item 将这棵二叉树转换成对应的树（或森林）。
	      \end{enumerate}

	\item 一棵二叉树的先序、中序、后序序列如下，其中部分未标出，请构造出该二叉树。
	      \begin{itemize}
		      \item 先序序列：\_\_ C D E \_\_ G H I \_\_ K
		      \item 中序序列：C B \_\_ F A \_\_ J K I G
		      \item 后序序列：\_ E F D B \_\_ J
	      \end{itemize}

	\item 设二叉树的存储结构如下：
	      \begin{center}
		      \begin{tabular}{|c|c|c|c|c|c|c|c|c|c|c|}
			      \hline
			             & 1 & 2 & 3 & 4 & 5 & 6 & 7 & 8 & 9 & 10 \\ \hline
			      Lchild & 0 & 0 & 0 & 9 & 4 & 0 & 0 & 0 & 0 & 0  \\ \hline
			      Data   & J & I & H & A & G & E & D & B & C & F  \\ \hline
			      Rchild & 0 & 0 & 0 & 0 & 0 & 0 & 0 & 0 & 0 & 0  \\ \hline
		      \end{tabular}
	      \end{center}
	      其中 BT 为树根结点的指针，其值为 6，Lchild, Rchild 分别为结点的左、右孩子指针域，Data 为结点的数据域。试完成下列各题：
	      \begin{enumerate}
		      \item 画出二叉树 BT 的逻辑结构；
		      \item 写出按前序、中序、后序遍历该二叉树所得到的结点序列；
		      \item 画出二叉树的中序线索树。
	      \end{enumerate}

	\item 设有正文 AADBAACACCDACACAAD，字符集为 A, B, C, D，设计一套二进制编码，使得上述正文的编码最短。

	\item 试构造一棵二叉树，包含权为 1, 4, 9, 16, 25, 36, 49, 64, 81, 100 等 10 个终端结点，且具有最小的加权路径长度 WPL。

	\item 以数据集 $\{2, 5, 7, 9, 13\}$ 为权值构造一棵哈夫曼树，并计算其带权路径长度。（请按左子树根结点的权小于等于右子树根结点的权的次序构造）

	\item 树和二叉树之间有什么样的区别与联系？
\end{enumerate}

\section*{算法题}

\begin{enumerate}
	\item 二叉树用链式方式存储，编写对二叉树进行先序，中序，后序遍历算法。
	\item 二叉树用链式方式存储，设计一个按层次顺序遍历二叉树的算法。
	\item 设计一个算法，判断一棵二叉树是否为完全二叉树。
	\item 计算一棵二叉树中的叶子结点数。
	\item 计算一棵二叉树中的单分支结点数。
	\item 求一棵二叉树的高度。
	\item 一棵二叉树，指针 t 指向根结点，求该二叉树中结点 n 的双亲结点。（若 n 无双亲，输出相应信息，若有，输出其双亲的值）
	\item 编写一个将二叉树中每个结点的左右孩子互换的算法。
	\item 写出在二叉树中查找值为 x 的结点的算法。
\end{enumerate}

\end{document}
